% ch2_d (书页 85-92 / p100-p107)
%--C-TAIL--
% ---- 书页 85 / p100 ----
外微分算符作用在 $p$-形式 $\omega$ 与 $q$-形式 $\eta$ 的乘积上时，服从一种修正版的莱布尼茨（Leibniz）法则：
\begin{equation*}
\mathrm{d}(\omega\wedge\eta)=(\mathrm{d}\omega)\wedge\eta+(-1)^{p}\,\omega\wedge(\mathrm{d}\eta). \tag{2.78}
\end{equation*}
我们鼓励你们自己证明这一点。

外微分值得特别关注的原因在于，它\emph{是一个张量}，即使在弯曲时空中也是如此；这一点与它的近亲——偏导数——不同。对于 $p=1$，我们可以从 1-形式的偏导数的变换律，即式 (2.38)，看出这一点；其中那个碍事的非张量项可以写成
\begin{equation*}
W_{\nu}\frac{\partial x^{\mu}}{\partial x^{\mu'}}\frac{\partial}{\partial x^{\mu}}\frac{\partial x^{\nu}}{\partial x^{\nu'}}=W_{\nu}\frac{\partial^{2}x^{\nu}}{\partial x^{\mu'}\partial x^{\nu'}}. \tag{2.79}
\end{equation*}
由于偏导数相互对易，这个表达式在 $\mu'$ 与 $\nu'$ 下是对称的。而外微分被定义为反对称化的偏导数，所以这一项会消失（对称表达式的反对称部分为零）。于是剩下的就是正确的张量变换律；推广到任意 $p$ 也很直接。所以外微分是一个堂堂正正的张量算符；然而，它并不能完全替代偏导数，因为它只对形式有定义。下一章我们将定义协变导数，它更接近我们通常所设想的“偏导数向任意流形的推广”。

关于外微分运算，另一个有趣的事实是：对任何形式 $A$，
\begin{equation*}
\mathrm{d}(\mathrm{d}A)=0, \tag{2.80}
\end{equation*}
它也常写作 $\mathrm{d}^{2}=0$。这一恒等式是 $\mathrm{d}$ 的定义以及偏导数可对易这一事实的推论：作用在任何对象上都有 $\partial_{\alpha}\partial_{\beta}=\partial_{\beta}\partial_{\alpha}$。这把我们引向下面一段数学题外话，纯粹为了好玩。我们把一个 $p$-形式 $A$ 定义为\textbf{闭}（closed）的，如果 $\mathrm{d}A=0$；称它是\textbf{恰当}（exact）的，如果 $A=\mathrm{d}B$，其中 $B$ 是某个 $(p-1)$-形式。显然，所有恰当形式都是闭形式，但反之则未必成立。在流形 $M$ 上，闭 $p$-形式构成一个矢量空间 $Z^{p}(M)$，恰当形式构成一个矢量空间 $B^{p}(M)$。我们定义一个新的矢量空间，其元素称为上同调类（cohomology class），即闭形式模掉恰当形式所得的空间：
\begin{equation*}
H^{p}(M)=\frac{Z^{p}(M)}{B^{p}(M)}. \tag{2.81}
\end{equation*}
也就是说，两个闭形式［$Z^{p}(M)$ 的元素］如果相差一个恰当形式［$B^{p}(M)$ 的元素］，它们就定义同一个上同调类［$H^{p}(M)$ 的元素］。奇妙的是，上同调空间 $H^{p}(M)$ 的维数只依赖于流形 $M$ 的拓扑。Minkowski 空间在拓扑上等价于 $\mathbf{R}^{4}$，后者乏味得很，于是当 $p>0$ 时所有的 $H^{p}(M)$ 都为零；而 $p=0$ 时我们有 $H^{0}(M)=\mathbf{R}$。因此在 Minkowski 空间中，除了零形式之外，所有闭形式都是恰当的；零形式不可能是恰当形式，因为根本不存在 $-1$-形式，可以让零形式成为其外微分。
% ---- 书页 86 / p101 ----
以这种方式抽取出关于拓扑的信息，着实令人惊讶，而这本质上涉及的是微分方程的解。

我们将要介绍的最后一种作用于微分形式的运算是\textbf{Hodge 对偶性}（Hodge duality）。我们在 $n$ 维流形上定义\textbf{Hodge 星算符}（Hodge star operator），它是一个从 $p$-形式到 $(n-p)$-形式的映射，
\begin{equation*}
(*A)_{\mu_{1}\cdots\mu_{n-p}}=\frac{1}{p!}\epsilon^{\nu_{1}\cdots\nu_{p}}{}_{\mu_{1}\cdots\mu_{n-p}}A_{\nu_{1}\cdots\nu_{p}}, \tag{2.82}
\end{equation*}
它把 $A$ 映为“$A$ 的对偶”。与我们其他作用于形式的运算不同，Hodge 对偶确实依赖于流形的度规［这一点应当显而易见，因为为了定义式 (2.82)，我们不得不提升了 Levi-Civita 张量的某些指标］。将 Hodge 星算符作用两次，得到的要么是原形式本身，要么是相差一个正负号的原形式：
\begin{equation*}
**A=(-1)^{s+p(n-p)}A, \tag{2.83}
\end{equation*}
其中 $s$ 是度规本征值中负号的个数。

关于 Hodge 对偶有两个事实。第一，Hodge 意义上的“对偶性”有别于矢量与对偶矢量之间的关系。“对偶性”的观念是指从一个空间到另一个空间的变换，其性质是把该变换做两次就回到原来的空间。矢量与 1-形式之间的对偶，以及 $p$-形式与 $(n-p)$-形式之间的 Hodge 对偶，对二者而言这一点都成立，这应当是清楚的。矢量空间之间对偶性的一个要求是：原来的空间与变换后的空间具有相同的维数；$p$-形式空间与 $(n-p)$-形式空间正是如此。

第二个事实涉及三维欧几里得空间中的微分形式。两个 1-形式的楔积的 Hodge 对偶给出另一个 1-形式：
\begin{equation*}
*(U\wedge V)_{i}=\epsilon_{i}{}^{jk}U_{j}V_{k}. \tag{2.84}
\end{equation*}
（所有前面的系数都恰好相消。）由于欧几里得空间中的 1-形式就像矢量一样，我们得到了一个从两个矢量到单个矢量的映射。你们应当让自己确信，这正是通常的叉乘（cross product）；而且 Levi-Civita 张量的出现解释了为什么叉乘在宇称变换下变号（两个坐标互换，等价地，基矢互换）。这正是叉乘只存在于三维的原因——因为只有在三维，我们才有一个从两个对偶矢量到第三个对偶矢量的有趣映射。

电动力学为微分形式的应用提供了一个格外有说服力的例子。从外微分的定义可以清楚地看到，方程 (1.97) 可以被简洁地表述为 2-形式 $F_{\mu\nu}$ 的闭性：
\begin{equation*}
\mathrm{d}F=0. \tag{2.85}
\end{equation*}
这是否意味着 $F$ 也是恰当的？是的；正如我们已指出的，Minkowski 空间拓扑平凡，所以所有闭形式都是恰当的。因此必定存在某个 1-形式 $A_{\mu}$，使得
\begin{equation*}
F=\mathrm{d}A. \tag{2.86}
\end{equation*}
% ---- 书页 87 / p102 ----
这个 1-形式就是大家熟悉的电磁学中的\textbf{矢势}（vector potential），其 0 分量由标势给出，$A_{0}=\Phi$，正如我们在第 1 章中讨论过的那样。规范不变性（gauge invariance）表现为：理论在 $A\to A+\mathrm{d}\lambda$ 下不变，其中 $\lambda$ 是某个标量（零形式）；从关系式 (2.86) 出发，这一点也是显而易见的。麦克斯韦方程组中的另一条，即 (1.96)，则可以表述为 3-形式之间的一个方程：
\begin{equation*}
\mathrm{d}(*F)=*J, \tag{2.87}
\end{equation*}
其中电流 1-形式 $J$ 不过是降低了指标的电流四维矢量。把细节补全就留给你们了，这是把微分形式记号转换为普通指标记号的绝佳练习。

Hodge 对偶性与某些场论的一个迷人性质密切相关：强耦合与弱耦合之间的对偶。很难不注意到，方程 (2.85) 与 (2.87) 看起来非常相似。的确，如果令 $J_{\mu}=0$，方程在如下“对偶变换”下不变：
\begin{gather*}
F\to *F,\\
*F\to -F. \tag{2.88}
\end{gather*}
因此我们说，真空中的麦克斯韦方程组是对偶不变的，而电荷的存在会破坏这种不变性。我们可以设想，自然界中除了电单极子之外还存在磁单极子；那样我们就可以在 (2.85) 的右端加上一个磁流项 $*J_{M}$，方程就会在下述组合下不变：对偶变换再加上一个额外的替换 $J\leftrightarrow J_{M}$。（当然，(2.85) 的右端若非零就与 $F=\mathrm{d}A$ 不相容，所以这个想法只有在 $A_{\mu}$ 不是基本变量时才行得通。）Dirac 考虑过磁单极子的想法，并证明了它们存在的一个必要条件：基本磁单极荷必须与基本电荷成反比。如今，基本电荷是一个小数目；电动力学是\emph{弱耦合}（weakly coupled）的，这正是微扰论在量子电动力学（QED）中如此惊人成功的原因。但 Dirac 关于磁荷的条件意味着，一个对偶变换会把一个弱耦合电荷的理论变成一个强耦合磁单极子的理论（反之亦然）。遗憾的是，磁单极子不太容易纳入普通电磁学，所以这些想法并不能直接应用；不过某种对偶对称性可能存在于某些理论之中（比如超对称非阿贝尔规范理论）。如果它确实存在，我们就有机会通过考察弱耦合的对偶版本，来分析一个看似强耦合（因而难以求解）的理论；在某些理论中，实际发生的正是这种情形。我们希望这些技术将使我们能够探索强耦合量子场论中确然存在的各种现象，比如夸克在强子中的禁闭。
% ---- 书页 88 / p103 ----
\section{积分}

张量密度与微分形式的一个重要应用场合是流形上的积分。你们大概已经接触过这样的事实：在 $\mathbf{R}^{n}$ 上的普通微积分中，体积元 $d^{n}x$ 在坐标变换下会多出一个雅可比行列式因子：
\begin{equation*}
d^{n}x'=\left|\frac{\partial x^{\mu'}}{\partial x^{\mu}}\right|d^{n}x. \tag{2.89}
\end{equation*}
其实，从微分形式的观点出发，这个公式有一个漂亮的解释，它来自如下事实：\emph{在 $n$ 维流形 $M$ 上，被积函数应当被恰当地理解为 $n$-形式。}换句话说，$n$ 维区域 $\Sigma\subset M$ 上的积分是一个从 $n$-形式场 $\omega$ 到实数的映射：
\begin{equation*}
\int_{\Sigma}:\omega\to\mathbf{R}. \tag{2.90}
\end{equation*}
这样的说法初听起来可能有些奇怪，但在曲线积分的语境里它肯定显得眼熟。在一维情形，任何一个 1-形式都可以写成 $\omega=\omega(x)\mathrm{d}x$，其中第一个 $\omega$ 是 1-形式，$\omega(x)$ 表示（唯一的那一个）分量函数。确实，我们把一维积分写成 $\int\omega(x)\mathrm{d}x$；你们也许习惯于把符号 $\mathrm{d}x$ 想成一个无穷小距离，但更确切地说，它是一个微分形式。

为了把这一点弄得更清楚，考虑不止一维的情形。如果我们宣称被积对象是一个 $n$-形式，就需要解释它在什么意义上是反对称的，而且顺便还要解释它为什么根本就是一个 $(0,n)$ 张量（从 $n$ 个矢量构成的集合到 $\mathbf{R}$ 的线性映射）。我们都同意积分可以写成 $\int f(x)\,d\mu$，其中 $f(x)$ 是流形上的标量函数，$d\mu$ 是体积元，或者叫测度。体积元的作用是给每个（无穷小的）区域指定一个（无穷小的）实数，即该区域的体积。无穷小区域（相对于有限大小的区域）的一个好处是，它们可以取成矩形的平行六面体——在存在曲率时，我们对“矩形平行六面体”应当是什么意思并没有清楚的概念，但当我们在无穷小区域中工作时，曲率的效应可以忽略。显然我们在这里并不严谨，但我们当下的目的纯粹是建立动机。

%FIG{2.27}: {一个无穷小的 $n$ 维区域，表示为一个平行六面体，由 $n$ 个矢量构成的有序集合定义，图中记为 $U$、$V$、$W$。}
如图 2.27 所示（为了便于说明，我们把流形取成三维的），一个平行六面体由定义其棱边的 $n$ 个矢量确定。于是，我们的体积元应当是一个从 $n$ 个矢量到实数的映射：$d\mu(U,V,W)\in\mathbf{R}$。（严格说来，它应当是一个从无穷小矢量到无穷小数的映射，但这样的映射同样会把有限矢量映到有限的数。）同样清楚的是，它应当可以按实数线性缩放；如果我们改变任何一个定义矢量的长度，体积就会相应地改变：$d\mu(aU,bV,cW)=abc\,d\mu(U,V,W)$。至于对矢量相加的线性性，就不是那么显而易见了，不过你们可以通过画图让自己信服。
% ---- 书页 89 / p104 ----
因此，我们的体积元是一个名副其实的 $(0,n)$ 张量。为什么反对称？因为我们定义的是一个有向的（oriented）对象；如果交换其中两个矢量，得到的体积应当大小相同而符号相反。（如果这一点不明显，你们至少应当让自己相信：当两个矢量共线时，体积应当为零。）这样，$n$ 维情形下的体积元在非常切实的意义上就是 $n$-形式。

要做实际计算，我们需要把这些想法变得更具体，而这其实很简单。关键的洞察是：把朴素的体积元 $d^{n}x$ 等同于用楔积构造出来的一个反对称张量密度：
\begin{equation*}
d^{n}x=\mathrm{d}x^{0}\wedge\cdots\wedge\mathrm{d}x^{n-1}. \tag{2.91}
\end{equation*}
右边的表达式可能有误导性，因为它看上去像一个张量（准确说是一个 $n$-形式），但其实是一个密度。当然，如果我们在 $M$ 上有两个函数 $f$ 和 $g$，那么 $\mathrm{d}f$ 和 $\mathrm{d}g$ 是 1-形式，而 $\mathrm{d}f\wedge\mathrm{d}g$ 是 2-形式。但出现在 (2.91) 中的函数是坐标函数本身，所以当我们变换坐标时，是把 1-形式 $\mathrm{d}x^{\mu}$ 换成一组新的 $\mathrm{d}x^{\mu'}$。你们看这里滑稽的地方——通常坐标变换改变的是分量，而不是 1-形式本身。(2.91) 的右端是一个依赖于坐标的对象（精确地说，是一个张量密度），它在 $x^{\mu}$ 坐标系中的行为就如同 $\mathrm{d}x^{0}\wedge\cdots\wedge\mathrm{d}x^{n-1}$。让我们看看它的实际运作。首先注意到，楔积的定义允许我们写出
\begin{equation*}
\mathrm{d}x^{0}\wedge\cdots\wedge\mathrm{d}x^{n-1}=\frac{1}{n!}\tilde{\epsilon}_{\mu_{1}\cdots\mu_{n}}\mathrm{d}x^{\mu_{1}}\wedge\cdots\wedge\mathrm{d}x^{\mu_{n}}, \tag{2.92}
\end{equation*}
这是因为楔积与 Levi-Civita 符号都是完全反对称的。（$1/n!$ 这个因子负责处理对指标的排列求和所引入的重复计数。）在坐标变换下，$\tilde{\epsilon}_{\mu_{1}\cdots\mu_{n}}$ 保持不变，而各个 1-形式按照式 (2.26) 变换，这就导致
\begin{align*}
\tilde{\epsilon}_{\mu_{1}\cdots\mu_{n}}\mathrm{d}x^{\mu_{1}}\wedge\cdots\wedge\mathrm{d}x^{\mu_{n}}&=\tilde{\epsilon}_{\mu_{1}\cdots\mu_{n}}\frac{\partial x^{\mu_{1}}}{\partial x^{\mu_{1}'}}\cdots\frac{\partial x^{\mu_{n}}}{\partial x^{\mu_{n}'}}\mathrm{d}x^{\mu_{1}'}\wedge\cdots\wedge\mathrm{d}x^{\mu_{n}'}\\
&=\left|\frac{\partial x^{\mu}}{\partial x^{\mu'}}\right|\tilde{\epsilon}_{\mu_{1}'\cdots\mu_{n}'}\mathrm{d}x^{\mu_{1}'}\wedge\cdots\wedge\mathrm{d}x^{\mu_{n}'}. \tag{2.93}
\end{align*}
在两边同乘雅可比行列式，并利用 (2.91) 与 (2.92)，便回到了 (2.89)。

显然，朴素的体积元 $d^{n}x$ 是按密度而非按张量变换的；不过，只要乘上 $\sqrt{|g|}$，就能直接构造出一个不变的体积元：
\begin{equation*}
\sqrt{|g'|}\,\mathrm{d}x^{0'}\wedge\cdots\wedge\mathrm{d}x^{(n-1)'}=\sqrt{|g|}\,\mathrm{d}x^{0}\wedge\cdots\wedge\mathrm{d}x^{n-1}, \tag{2.94}
\end{equation*}
它当然就正是 $(n!)^{-1}\epsilon_{\mu_{1}\cdots\mu_{n}}\mathrm{d}x^{\mu_{1}}\wedge\cdots\wedge\mathrm{d}x^{\mu_{n}}$。为简单起见，我们通常把体积元写成 $\sqrt{|g|}\,d^{n}x$，而不写成显式的楔积：
% ---- 书页 90 / p105 ----
\begin{equation*}
\sqrt{|g|}\,d^{n}x\equiv\sqrt{|g|}\,\mathrm{d}x^{0}\wedge\cdots\wedge\mathrm{d}x^{n-1}\ ; \tag{2.95}
\end{equation*}
心里记住它本应是一个 $n$-形式，就足够了。事实上，体积元不多不少，恰恰就是 Levi-Civita 张量 $\epsilon_{\mu_{1}\cdots\mu_{n}}$；把基 1-形式显式地补写回来，我们看到
\begin{align*}
\epsilon&\equiv\epsilon_{\mu_{1}\cdots\mu_{n}}\mathrm{d}x^{\mu_{1}}\otimes\cdots\otimes\mathrm{d}x^{\mu_{n}}\\
&=\frac{1}{n!}\epsilon_{\mu_{1}\cdots\mu_{n}}\mathrm{d}x^{\mu_{1}}\wedge\cdots\wedge\mathrm{d}x^{\mu_{n}}\\
&=\frac{1}{n!}\sqrt{|g|}\,\tilde{\epsilon}_{\mu_{1}\cdots\mu_{n}}\mathrm{d}x^{\mu_{1}}\wedge\cdots\wedge\mathrm{d}x^{\mu_{n}}\\
&=\sqrt{|g|}\,\mathrm{d}x^{0}\wedge\cdots\wedge\mathrm{d}x^{n-1}\\
&\equiv\sqrt{|g|}\,d^{n}x. \tag{2.96}
\end{align*}
注意到，epsilon 张量引入的组合因子，恰好与从张量积换成楔积时引入的因子相消；之所以允许做这种替换，是因为 epsilon 张量会自动反对称化。

于是，结论很简单：标量函数 $\phi$ 在 $n$ 维流形上的积分 $I$ 写作
\begin{equation*}
\boxed{I=\int\phi(x)\sqrt{|g|}\,d^{n}x.} \tag{2.97}
\end{equation*}
给定 $\phi(x)$ 和 $\sqrt{|g|}$ 的显式表达式，这样的积分就可以用多变量微积分的通常方法直接算出。度规的行列式所起的作用是自动保证正确的变换性质。你们有时会看到更抽象的记号
\begin{equation*}
I=\int\phi(x)\,\epsilon\ ; \tag{2.98}
\end{equation*}
鉴于 (2.96)，这两种写法传递的是同样的内容。

\section{习题}

% 习题编号照原书
\textbf{1.}\quad 一个流形在拓扑上非平凡，并不必然意味着它不能用单个卡片覆盖。与圆周 $S^{1}$ 的情形相对照，请通过显式地构造映射，证明无穷长圆柱 $\mathbf{R}\times S^{1}$ 可以只用一个卡片覆盖。

\textbf{2.}\quad 通过巧妙地选取坐标卡片，我们能让 $\mathbf{R}^{2}$ 看起来像一个一维流形吗？能让 $\mathbf{R}^{1}$ 看起来像一个二维流形吗？如果能，请显式构造一个适当的图册；如果不能，请解释为什么不能。这个问题的要点在于促使你深入思考流形究竟是什么；若不进一步讨论拓扑空间的更多细节，就无法严格地回答它。特别是，你可能必须忘掉你在原来的 $\mathbf{R}^{2}$ 或 $\mathbf{R}^{1}$ 中已经知道的那套“开集”定义，而把它们定义为恰当地继承自所映射到的 $\mathbf{R}^{1}$ 或 $\mathbf{R}^{2}$。

% ---- 书页 91 / p106 ----
\textbf{3.}\quad 请通过显式构造一个适当的图册，证明二维环面 $T^{2}$ 是一个流形。（显然，不必是极大的图册。）

\textbf{4.}\quad 验证 2.3 节末尾关于两个矢量场的对易子所作的那些论断（线性性、Leibniz 律、分量公式、作为一个矢量场的变换）。

\textbf{5.}\quad 在 $\mathbf{R}^{2}$ 中给出两个线性无关、处处不为零的矢量场的例子，使它们的对易子不为零。注意，这两个场在每一点都给出切空间的一组基，但由于对易子不为零，它不可能是坐标基。

\textbf{6.}\quad 把 $\mathbf{R}^{3}$ 看作一个带有平直欧几里得度规的流形，坐标为 $\{x,y,z\}$。引入球坐标 $\{r,\theta,\phi\}$，其与 $\{x,y,z\}$ 的关系为
\begin{gather*}
x=r\sin\theta\cos\phi\\
y=r\sin\theta\sin\phi\\
z=r\cos\theta, \tag{2.99}
\end{gather*}
于是度规取如下形式
\begin{equation*}
ds^{2}=\mathrm{d}r^{2}+r^{2}\mathrm{d}\theta^{2}+r^{2}\sin^{2}\theta\,\mathrm{d}\phi^{2}. \tag{2.100}
\end{equation*}

\textbf{(a)}\quad 一个粒子沿如下参数化曲线运动：
\begin{equation*}
x(\lambda)=\cos\lambda,\qquad y(\lambda)=\sin\lambda,\qquad z(\lambda)=\lambda. \tag{2.101}
\end{equation*}
请在 $\{r,\theta,\phi\}$ 系中表出这条曲线的路径。

\textbf{(b)}\quad 分别在笛卡尔坐标系与球坐标系中，计算该曲线的切矢量的分量。

\textbf{7.}\quad 长椭球坐标（prolate spheroidal coordinates）可以用来简化天体力学中的 Kepler 问题。它们与欧几里得三维空间中通常的笛卡尔坐标 $(x,y,z)$ 的关系为
\begin{gather*}
x=\sinh\chi\,\sin\theta\cos\phi,\\
y=\sinh\chi\,\sin\theta\sin\phi,\\
z=\cosh\chi\,\cos\theta.
\end{gather*}
把注意力限制在 $y=0$ 的平面上，回答下列问题。

\textbf{(a)}\quad 联系 $(x,z)$ 与 $(\chi,\theta)$ 的坐标变换矩阵 $\partial x^{\mu}/\partial x^{\nu'}$ 是什么？

\textbf{(b)}\quad 线元 $ds^{2}$ 在长椭球坐标下是什么样子？

\textbf{8.}\quad 验证式 (2.78)：对 $p$-形式 $\omega$ 与 $q$-形式 $\eta$ 的乘积的外微分，我们有
\begin{equation*}
\mathrm{d}(\omega\wedge\eta)=(\mathrm{d}\omega)\wedge\eta+(-1)^{p}\,\omega\wedge(\mathrm{d}\eta). \tag{2.102}
\end{equation*}

% ---- 书页 92 / p107 ----
\textbf{9.}\quad 在 Minkowski 空间中，设 $*F=q\sin\theta\,\mathrm{d}\theta\wedge\mathrm{d}\phi$。

\textbf{(a)}\quad 计算 $\mathrm{d}*F=*J$。

\textbf{(b)}\quad 2-形式 $F$ 等于什么？

\textbf{(c)}\quad 对这个解来说，电场与磁场分别等于什么？

\textbf{(d)}\quad 计算 $\int_{V}\mathrm{d}*F$，其中 $V$ 是在某个固定时刻、欧几里得三维空间中半径为 $R$ 的一个球。

\textbf{10.}\quad 考虑二维时空中的麦克斯韦方程组 $\mathrm{d}F=0$、$\mathrm{d}*F=*J$。解释为什么这两组方程中可以舍弃一组。证明电磁场可以用一个标量场表出，并写出这个标量场的场方程的分量形式。

\textbf{11.}\quad 本题为一个非常简单的问题附加了大量动机性的文字；别因此分心。在与点粒子相联系的普通电磁学中，作用量中代表规范势 1-形式 $A^{(1)}$ 与带电粒子相耦合的部分可以写成 $S=\int_{\gamma}A^{(1)}$，其中 $\gamma$ 是粒子的世界线。（$A^{(1)}$ 的上标只是为了提醒你它是一个 1-形式。）本题中你要考虑一个与普通电磁学有亲缘关系的理论，但这一次是在 11 维时空里，带有 3-形式规范势 $A^{(3)}$ 和 4-形式场强 $F^{(4)}=\mathrm{d}A^{(3)}$。注意，场强在规范变换 $A^{(3)}\to A^{(3)}+\mathrm{d}\lambda^{(2)}$（$\lambda^{(2)}$ 为任意 2-形式）下不变。

\textbf{(a)}\quad 这种规范场自然与之耦合的对象，其空间维数应当是多少？（例如，普通电磁学（E+M）与零维对象——点粒子——相耦合。）

\textbf{(b)}\quad 普通电子的电荷由如下积分给出：2-形式规范场强的对偶在围绕粒子的一个二维球面上的积分。对于 $A^{(3)}$ 所耦合的对象，你会如何定义它的“电荷”？请论证：如果 $\mathrm{d}*F^{(4)}=0$，这个荷是守恒的。

\textbf{(c)}\quad 设想存在一个“对偶规范势” $\widetilde{A}$，满足 $\mathrm{d}(\widetilde{A})=*F^{(4)}$。它自然与之耦合的对象是多少维的？

\textbf{(d)}\quad 规范场自身的作用量（相对于它与其他对象的耦合而言）将是整个 11 维时空上的一个积分。这样的作用量中，容许出现而且在“局域”规范变换下不变的项有哪些？例如，规范变换可以由一个在无穷远处为零的 2-形式 $\lambda^{(2)}$ 来规定。请把自己限制在关于 $A^{(3)}$ 及其一阶导数的线性、二次或三次的项（不含二阶导数，不含更高阶的项）。你可以使用外微分、楔积和 Hodge 对偶，但不允许度规的任何显式出现。

更多的背景知识：“超对称性”（supersymmetry）是一种假想中的对称性，它把玻色子（自旋为整数的粒子）与费米子（自旋为 $\frac{1}{2}$、$\frac{3}{2}$ 等的粒子）联系起来。一个有趣的性质是，超对称理论只有在 11 维或更低维数中才是良定义的——在更高的维数下，超对称性将要求存在自旋大于 2 的粒子，而这样的粒子无法自洽地量子化。十一维超对称是一个独一无二的理论，它自然地包含一个 3-形式规范势（更不用说引力了）。近期的工作表明，它还包含本题所暗指的各种更高维的对象（尽管我们在这里走了些捷径）。这个理论最终被发现是某个称作 $M$-理论（M-theory）的东西的一个良定义的极限，而 $M$-理论的其他极限则是各种 10 维超弦理论。
